\section{Synthesis of the multi-paradigm approach}
Having detailed our methodological journey, we now turn to the empirical evaluation of our systems. To comprehensively address the challenges of the SemEval-2025 narrative classification task, we developed a suite of models spanning three distinct paradigms. We began with \textbf{Architecture 0}, a traditional \textbf{fine-tuning} approach based on \textbf{mDeBERTa}, systematically enhancing it with techniques like \textbf{class weighting} and \textbf{hierarchical regularization} to establish a robust \textbf{supervised baseline}. Recognizing the limitations of fine-tuning on a \textbf{data-scarce, long-tailed} problem, we then shifted to a \textbf{zero-shot reasoning} paradigm. This led to \textbf{Architecture 1}, a novel \textbf{agentic framework} that decomposes the task for specialized \textbf{LLM agents}, and ultimately to \textbf{Architecture 2}, our most advanced \textbf{traceable actor-critic pipeline}, designed for maximum \textbf{reliability} and \textbf{interpretability} through \textbf{evidence-grounded validation}. In this section, we present a comparative analysis of these three paradigms to evaluate their effectiveness and trade-offs.